\corrigesChapitre{5}{Atomes, molécules et formules}
\par\noindent\begin{minipage}{\linewidth}
\begin{corrige}{1}{Quatre symboles}
\begin{questions}
\item H et O.
\item C : carbone ; N : azote. Les majuscules sont indispensables.
\end{questions}
\end{corrige}
\end{minipage}
\par\noindent\begin{minipage}{\linewidth}
\begin{corrige}{2}{Lire l’eau}
\begin{questions}
\item Deux H et un O.
\item Trois atomes par molécule.
\end{questions}
\end{corrige}
\end{minipage}
\par\noindent\begin{minipage}{\linewidth}
\begin{corrige}{3}{Deux molécules différentes}
\begin{questions}
\item $\ce{CO}$ contient un C et un O ; $\ce{CO2}$, un C et deux O.
\item Non. Les propriétés des deux espèces sont différentes ; le monoxyde de carbone est notamment très toxique.
\end{questions}
\end{corrige}
\end{minipage}
\par\noindent\begin{minipage}{\linewidth}
\begin{corrige}{4}{Retrouver un nom}
\begin{questions}
\item Diazote et dihydrogène.
\item $\ce{CH4}$.
\end{questions}
\end{corrige}
\end{minipage}
\par\noindent\begin{minipage}{\linewidth}
\begin{corrige}{5}{Repérer une charge}
\begin{questions}
\item $\ce{Na+}$ et $\ce{Cl-}$.
\item Le premier porte une charge positive ; le second une charge négative. La molécule d’eau est neutre.
\end{questions}
\end{corrige}
\end{minipage}
\par\noindent\begin{minipage}{\linewidth}
\begin{corrige}{6}{Plusieurs molécules}
\begin{questions}
\item $5$ atomes C et $10$ atomes O.
\item $5+10=15$ atomes, répartis en cinq molécules.
\end{questions}
\par\smallskip\noindent\textbf{Pour comprendre.}\enspace Le coefficient 5 multiplie toute la formule CO2 : $5\times1=5$ atomes C et $5\times2=10$ atomes O, donc 15 atomes au total.\par
\end{corrige}
\end{minipage}
\par\noindent\begin{minipage}{\linewidth}
\begin{corrige}{7}{Le protoxyde d’azote}
\begin{questions}
\item Deux N et un O.
\item Six N et trois O, soit neuf atomes au total. Le préfixe du nom ne remplace pas la lecture de la formule.
\end{questions}
\par\smallskip\noindent\textbf{Pour comprendre.}\enspace Une molécule N2O contient deux N et un O. Trois molécules contiennent $3\times2=6$ N et $3\times1=3$ O. L’indice ne doit pas être confondu avec le coefficient.\par
\end{corrige}
\end{minipage}
\par\noindent\begin{minipage}{\linewidth}
\begin{corrige}{8}{Une case du tableau}
\begin{questions}
\item Al est le symbole ; 13 est le numéro atomique.
\item Un symbole à deux lettres comporte une majuscule puis une minuscule.
\end{questions}
\end{corrige}
\end{minipage}
\par\noindent\begin{minipage}{\linewidth}
\begin{corrige}{9}{Un mélange de modèles}
\begin{questions}
\item Deux espèces : dioxygène et diazote.
\item Huit O et quatre N. Le nombre de molécules est six, différent du nombre d’atomes.
\end{questions}
\pcFigureEx{figures/cor-melange-molecules.pdf}{Chaque boule est un atome ; les boules liées forment une molécule. Compter séparément les deux espèces.}
\end{corrige}
\end{minipage}
\par\noindent\begin{minipage}{\linewidth}
\begin{corrige}{10}{Corriger une formule}
\begin{questions}
\item Il faut écrire $\ce{H2O}$.
\item $\ce{HO2}$ décrirait une composition avec un H et deux O ; ce n’est pas la molécule d’eau.
\end{questions}
\pcFigureEx{figures/cor-eau-formule.pdf}{L’indice 2 doit porter sur H : il y a deux H et un O.}
\end{corrige}
\end{minipage}
\par\noindent\begin{minipage}{\linewidth}
\begin{corrige}{11}{Même inventaire, autre assemblage}
\begin{questions}
\item Dans chaque cas, quatre H et deux O.
\item À gauche ils sont assemblés en eau ; à droite en dihydrogène et dioxygène. L’inventaire des atomes ne suffit pas à déterminer l’assemblage.
\end{questions}
\pcFigureEx{figures/cor-inventaires.pdf}{Même inventaire de quatre H et deux O ; assemblages et espèces différents.}
\end{corrige}
\end{minipage}
\par\noindent\begin{minipage}{\linewidth}
\begin{corrige}{12}{Une collection à construire}
\begin{questions}
\item Chaque modèle demande quatre H ; $20/4=5$ modèles sont possibles.
\item Cinq C et vingt H ont été utilisés ; il reste un C. On ne modifie pas la formule pour utiliser tous les jetons.
\end{questions}
\par\smallskip\noindent\textbf{Pour comprendre.}\enspace Chaque modèle CH4 demande un C et quatre H. Les maxima sont $6/1=6$ et $20/4=5$ ; retenir le plus petit. Cinq modèles utilisent 5 C et 20 H : il reste un C.\par
\pcFigureEx{figures/cor-methane-jetons.pdf}{Le carbone permettrait six modèles ; l’hydrogène n’en permet que cinq : il limite la construction.}
\end{corrige}
\end{minipage}
\par\noindent\begin{minipage}{\linewidth}
\begin{corrige}{13}{Une notation ambiguë}
\begin{questions}
\item Co est le symbole du cobalt ; CO est la formule du monoxyde de carbone.
\item Respecter les majuscules et minuscules, et nommer l’espèce ou l’élément à la première utilisation.
\end{questions}
\end{corrige}
\end{minipage}
\par\noindent\begin{minipage}{\linewidth}
\begin{corrige}{14}{Reconstituer une formule}
\begin{questions}
\item $\ce{C2H6O}$.
\item Non. La formule brute donne la composition, sans toujours déterminer l’organisation des atomes. Cet exemple annonce une limite du modèle.
\end{questions}
\end{corrige}
\end{minipage}
\par\noindent\begin{minipage}{\linewidth}
\begin{corrige}{15}{Critiquer une représentation}
\begin{questions}
\item Des atomes et molécules ne mesurent pas un millimètre. Leur taille a été énormément agrandie.
\item Annoncer une représentation schématique non à l’échelle ; la couleur code une espèce ou un élément selon la légende.
\end{questions}
\end{corrige}
\end{minipage}
\par\noindent\begin{minipage}{\linewidth}
\phantomsection\label{cor-defi-5}\pcSolutionDefi{5}{Une commande de modèles}\begin{questions}
\item H : $16$ ; O : $8+10=18$ ; C : $5$.
\item Oui : $8+5=13$ molécules et $16+18+5=39$ atomes dans le modèle.
\item Un comptage d’atomes n’est pas une mesure de masse ; les atomes n’ont pas tous la même masse et les jetons ne sont pas des atomes réels.
\end{questions}
\end{minipage}
\par\noindent\begin{minipage}{\linewidth}
\phantomsection\label{cor-activite-5}\pcSousTitre{Pistes pour l’investigation}
Vérifier pour chaque modèle le nombre de H, O, C ou N indiqué par les indices. Deux boules O reliées font une molécule de dioxygène ; deux molécules exigent quatre boules. Les couleurs sont un code. La construction ne montre ni les tailles réelles, ni les distances, ni les vibrations ; elle sert d’abord à représenter composition et assemblage.\par\end{minipage}
